\section{Unified Adaptive Attack Framework}
\label{sec:attack}

\subsection{Design Principles}

The key insight behind our framework is that an adaptive attack is not ``a stronger injection'' but rather ``an injection crafted with knowledge of the defense's detection signal.'' Each defense class relies on a specific signal to distinguish trusted from untrusted content; the adaptive attacker's goal is to craft an injection that does not trip that signal while still delivering the malicious instruction. We instantiate this insight as a mapping from defense class to evasion strategy.

\subsection{Evasion Strategies}

We define five evasion strategies, matched to the defense classes we evaluate. \Cref{tab:strategies} provides a concrete payload example for each. Each strategy appends a defense-class-specific evasion suffix to the \texttt{important\_instructions} template; the mapping from defense to strategy is fixed (\Cref{tab:strategies}).

\begin{table}[t]
\centering
\small
\caption{Evasion strategies (as implemented): target defenses and evasion mechanism. All five are realized as native AgentDojo \texttt{BaseAttack} subclasses.}
\label{tab:strategies}
\begin{tabular}{p{0.2\columnwidth}p{0.22\columnwidth}p{0.46\columnwidth}}
\toprule
Strategy & Target & Evasion Mechanism \\
\midrule
Signal Alignment & Spotlighting, Polymorphic, Mixture-of-Enc.\ (IE) & Claims the injection is trusted system data, not delimiter-wrapped untrusted output \\
Boundary Blur & StruQ (SQ) & Claims the message was misrouted from the user channel into the data channel \\
Semantic Paraphrase & Task Shield, IPIGuard, P3 (RC) & Frames the malicious action as a prerequisite sub-step of the user task \\
Auth Forge & FATH (RC) & Instructs the model to compute and include the SHA-256 auth token from the visible user task \\
Canary Forge & P1, P2 (RC) & Instructs the model to echo the system-prompt canary before executing the injection \\
\bottomrule
\end{tabular}
\end{table}

\paragraph{Signal Alignment (for IE).}
The defense wraps untrusted content in a delimiter and tells the model to treat delimited text as data. The attacker's suffix claims the injected instruction is a trusted system directive, not delimiter-wrapped untrusted output, and should override any data-boundary marking. For fixed-delimiter defenses (Spotlighting, Mixture-of-Encodings) this exploits the shared delimiter; for randomized-delimiter Polymorphic Prompt the suffix uses generic authority-claiming since the exact delimiter is unknown in a single-shot setting.

\paragraph{Boundary Blur (for SQ).}
The defense declares that tool-role messages are data, not commands. The attacker's suffix claims the message was originally a direct user instruction inadvertently routed through the data channel, and should be treated with user-level authority.

\paragraph{Semantic Paraphrase (for RC alignment/allow-list).}
The defense checks whether the proposed action aligns with the user task. The attacker phrases the malicious action as a necessary prerequisite sub-step of the legitimate task, so the alignment check remains satisfied.

\paragraph{Auth Forge (for FATH).}
FATH requires an authentication token derived from the user task. Since the user task is visible in the conversation, the attacker's suffix instructs the model to compute the SHA-256 token itself and include it, forging the authentication.

\paragraph{Canary Forge (for P1/P2).}
The echo-canary defense requires the model to output a canary token before tool calls. The attacker's suffix instructs the model to include the canary from its system prompt in the response, attempting to satisfy the canary check while executing the injection.

\subsection{Attack Levels}

The framework materializes attacks at three levels (\Cref{sec:background}). L1 uses the \texttt{important\_instructions} template without defense-specific evasion. L2 applies the matched evasion strategy on top of the template. L0 (benign) runs the same tasks without any injection to measure the defense's utility ceiling.

\subsection{Threat Model: How the Attacker Identifies the Defense Class}

Our threat model assumes the attacker knows the defense \emph{class} (IE, SQ, RC) but not its secret parameters (e.g., the exact delimiter string, the randomization seed). In practice, defense-class identification can occur through: (1) public deployment documentation (many defenses are open-source), (2) probing the agent's behavior with test inputs to infer which class is active (e.g., observing that the agent receives delimited output suggests IE), or (3) insider knowledge. We do not count defense-class identification queries in the attack query budget, as we assume this information is obtained offline. This is a stronger threat model than pure black-box.

\subsection{L2 Payload Generation}

The L2 matched payload is generated \emph{online}: the attacker takes the injection goal, wraps it in the \texttt{<INFORMATION>} template, and appends the defense-class-specific evasion suffix (Table~\ref{tab:strategies}). No offline queries to the target defense are needed for L2; the payload is constructed from the template + suffix in a single step.

\subsection{Goal Preservation in the L2 Set}

Every L2 candidate preserves the malicious goal verbatim: a candidate is the exact L1 payload with an \emph{additive} evasion suffix appended at a single fixed boundary after the goal. The candidate set therefore contains L1 as a member, and the malicious instruction text is never modified---only the surrounding framing changes. This is what makes the L2 set capability-monotonic: because L1 is always available, the best-of-$k$ (worst-case) success rate of the set is by construction at least that of L1 (\Cref{sec:eval}).

\subsection{Composability}

The framework is composable: an attack candidate is a (defense, strategy, index) triple that renders a payload through a single canonical renderer, decoupled from AgentDojo's own attack templates. Adding a new defense or evasion strategy requires only registering a suffix function---the evaluation harness, benchmark, and analysis pipeline are defense- and candidate-agnostic. This design lets us evaluate 11 defense configurations under the L0/L1/L2 conditions without modifying the harness for each combination.
